%% kapa: exam=1 round=37 year=2026 period=1 form=A subject=economics subject_ko=경제학원론 seq=27 number=67
%% source: Q-Net 감정평가사 기출문제 (cst00309) artlSeq=5258924 file=raw/1008/5258924/2026년 제37회 감정평가사 1차 1교시.pdf
%% conversion: kordoc 4.24.51 -> tools/gen1.py (auto); review=proofread-sonnet; flags=none
\begin{question}{67}
통화량 변동에 관한 설명으로 옳은 것을 모두 고른 것은?
\begin{framedtext}
ㄱ. 중앙은행의 외환보유고 증가는 통화량 감소를 가져온다. \par
ㄴ. 중앙은행의 국채매입은 통화량의 증가를 가져온다. \par
ㄷ. 시중은행의 자기자본비율(BIS비율)의 증가는 통화량 증가를 가져온다. \par
ㄹ. 시중은행의 초과지급준비율의 상승은 통화량 감소를 가져온다. \par
ㅁ. 재할인율의 인상은 통화량 감소를 가져온다. \par
\end{framedtext}
\begin{choices}
  \item ㄱ, ㄴ
  \item ㄱ, ㄷ, ㄹ
  \item ㄴ, ㄷ, ㅁ
  \item ㄴ, ㄹ, ㅁ
  \item ㄷ, ㄹ, ㅁ
\end{choices}
\end{question}
\answer{4}
